% Options for packages loaded elsewhere
\PassOptionsToPackage{unicode}{hyperref}
\PassOptionsToPackage{hyphens}{url}
\PassOptionsToPackage{dvipsnames,svgnames,x11names}{xcolor}
\documentclass[
  11pt,
  a4paper,
]{article}
\usepackage{xcolor}
\usepackage[margin=1in]{geometry}
\usepackage{amsmath,amssymb}
\setcounter{secnumdepth}{-\maxdimen} % remove section numbering
\usepackage{iftex}
\ifPDFTeX
  \usepackage[T1]{fontenc}
  \usepackage[utf8]{inputenc}
  \usepackage{textcomp} % provide euro and other symbols
\else % if luatex or xetex
  \usepackage{unicode-math} % this also loads fontspec
  \defaultfontfeatures{Scale=MatchLowercase}
  \defaultfontfeatures[\rmfamily]{Ligatures=TeX,Scale=1}
\fi
\usepackage{lmodern}
\ifPDFTeX\else
  % xetex/luatex font selection
\fi
% Use upquote if available, for straight quotes in verbatim environments
\IfFileExists{upquote.sty}{\usepackage{upquote}}{}
\IfFileExists{microtype.sty}{% use microtype if available
  \usepackage[]{microtype}
  \UseMicrotypeSet[protrusion]{basicmath} % disable protrusion for tt fonts
}{}
\makeatletter
\@ifundefined{KOMAClassName}{% if non-KOMA class
  \IfFileExists{parskip.sty}{%
    \usepackage{parskip}
  }{% else
    \setlength{\parindent}{0pt}
    \setlength{\parskip}{6pt plus 2pt minus 1pt}}
}{% if KOMA class
  \KOMAoptions{parskip=half}}
\makeatother
\setlength{\emergencystretch}{3em} % prevent overfull lines
\providecommand{\tightlist}{%
  \setlength{\itemsep}{0pt}\setlength{\parskip}{0pt}}
\usepackage{mathtools}
\usepackage{fancyhdr}

\allowdisplaybreaks
\emergencystretch=3em

\pagestyle{fancy}
\fancyhf{}
\fancyhead[L]{\small Erd\H{o}s Problem 625}
\fancyhead[R]{\small Proposed proof}
\fancyfoot[C]{\thepage}
\setlength{\headheight}{14pt}

\AtBeginDocument{%
  % Latin Modern Math lacks unicode-math's U+29F5 set-minus glyph.
  \renewcommand{\setminus}{\mathbin{\backslash}}%
}
\usepackage{bookmark}
\IfFileExists{xurl.sty}{\usepackage{xurl}}{} % add URL line breaks if available
\urlstyle{same}
\hypersetup{
  pdftitle={A self-contained proof of a polynomial-scale gap between the chromatic and cochromatic numbers of a random graph},
  colorlinks=true,
  linkcolor={blue},
  filecolor={Maroon},
  citecolor={Blue},
  urlcolor={blue},
  pdfcreator={LaTeX via pandoc}}

\title{A self-contained proof of a polynomial-scale gap between the
chromatic and cochromatic numbers of a random graph}
\usepackage{etoolbox}
\makeatletter
\providecommand{\subtitle}[1]{% add subtitle to \maketitle
  \apptocmd{\@title}{\par {\large #1 \par}}{}{}
}
\makeatother
\subtitle{Candidate all-n solution to Erd\H{o}s Problem 625; not
externally peer reviewed}
\author{}
\date{12 July 2026}

\begin{document}
\maketitle

{
\setcounter{tocdepth}{1}
\tableofcontents
}
\clearpage

\section*{Abstract}\label{abstract}
\addcontentsline{toc}{section}{Abstract}

For a graph \texttt{G}, let \texttt{chi(G)} be its chromatic number and
let \texttt{zeta(G)} be the least number of parts in a vertex partition
in which every part induces either an empty graph or a complete graph.
We prove that, for \texttt{G\_n\textasciitilde{}G(n,1/2)},

\[
 \Pr\left\{\chi(G_n)-\zeta(G_n)
 \ge \frac{(\ln2)^2}{32}\ln\frac{200}{153}
       \frac{n}{(\ln n)^3}\right\}\longrightarrow1.      \tag{0.1}
\]

The construction uses classes of the four sizes
\texttt{alpha-2,}\allowbreak{}\texttt{...,}\allowbreak{}\texttt{alpha-5},
where \texttt{alpha} is the usual independence-number phase. Signing
each class as independent or complete moves its first-moment root down
by order \texttt{n/(ln\ n)\^{}3}, uniformly through the whole phase
cycle. An exact sign-summed overlap identity reduces the second moment
to a bipartite configuration model. We sum partial diagonals,
unequal-type dense containments, all intermediate high cells, and every
residual even-subgraph attachment with logarithmic cost
\texttt{o(n/(ln\ n)\^{}4)}. Paley--Zygmund gives a possibly rare
cocolouring, and a proved Alon--Scott-type leftover argument amplifies
it to high probability without losing the root separation.

\section{1. Notation and elementary
facts}\label{notation-and-elementary-facts}

All logarithms are natural unless a base is displayed. Put

\[
 q=\ln2,\qquad N=\ln n,\qquad w=\ln N.                    \tag{1.1}
\]

For an integer \texttt{s}, let

\[
 \mu_s(v)=\binom vs2^{-\binom s2},\qquad \mu_s=\mu_s(n). \tag{1.2}
\]

We use the following standard inequalities in their stated forms.

\begin{enumerate}
\def\labelenumi{\arabic{enumi}.}
\tightlist
\item
  For integers \texttt{m\textgreater{}=1}, \[
   \ln(m!)=m\ln m-m+\frac12\ln(2\pi m)+O(1/m).          \tag{1.3}
  \] In particular, uniformly for \texttt{m\textgreater{}=0}, with
  \texttt{0\ ln0=0}, the last two terms may be replaced by
  \texttt{O(ln(m+1))}.
\item
  If a random variable \texttt{Y} is a function of \texttt{r}
  independent blocks and changing one block changes \texttt{Y} by at
  most one, then \[
   \Pr\{|Y-\mathbb EY|\ge t\}
   \le2\exp(-2t^2/r).                                    \tag{1.4}
  \] We will use the corresponding one-sided bounds as well.
\item
  If \texttt{Z\textgreater{}=0} and
  \texttt{0\textless{}EZ\textless{}infinity}, then \[
   \Pr\{Z>0\}\ge\frac{(\mathbb EZ)^2}{\mathbb EZ^2}.   \tag{1.5}
  \]
\item
  If \texttt{X\textasciitilde{}Bin(m,1/2)}, then \[
   \Pr\{X\le m/4\}\le e^{-m/16}.                       \tag{1.6}
  \] Indeed, exponential Markov with \texttt{t=ln\ 3} bounds the
  probability by
  \texttt{{[}3\^{}(1/}\allowbreak{}\texttt{4)}\allowbreak{}\texttt{(2/}\allowbreak{}\texttt{3)}\allowbreak{}\texttt{{]}}\allowbreak{}\texttt{\^{}m\textless{}}\allowbreak{}\texttt{=}\allowbreak{}\texttt{e\^{}\{-m/}\allowbreak{}\texttt{16\}}.
\item
  For every nonnegative random variable \texttt{X} and
  \texttt{a\textgreater{}0},
  \texttt{P(X\textgreater{}=a)\textless{}=E\ X/a}.
\end{enumerate}

The first is Stirling's estimate, the second is McDiarmid's bounded-
differences inequality, and the third is the zero-threshold case of
Paley--Zygmund. The fourth is the only binomial-tail estimate used
below; the fifth is Markov's inequality.

\section{2. The complete independence-number
phase}\label{the-complete-independence-number-phase}

Define

\[
 \alpha_0=2\log_2n-2\log_2\log_2n
             +2\log_2(e/2)+1,
 \quad \alpha=\lfloor\alpha_0\rfloor,
 \quad \delta=\alpha_0-\alpha.                           \tag{2.1}
\]

Thus \texttt{0\textless{}=delta\textless{}1}. The estimates below are
uniform on this entire interval, including sequences approaching either
endpoint.

\subsection{Lemma 2.1 (phase
expansion)}\label{lemma-2.1-phase-expansion}

There is a bounded continuous function \texttt{K(delta)} such that

\[
 \ln\mu_\alpha
 =\delta N+\left(\frac2q-\frac12-\delta\right)w
   +K(\delta)+O(w^2/N).                                  \tag{2.2}
\]

In particular,

\[
 \mu_{\alpha+2}=n^{\delta-2+o(1)}
 \le n^{-1+o(1)}=o(1),                                   \tag{2.3}
\]

and, with \texttt{a=alpha-2},

\[
 \mu_a\ge c n^2N^{2/q-5/2}                              \tag{2.4}
\]

for an absolute \texttt{c\textgreater{}0} and all sufficiently large
\texttt{n}.

\subsubsection{Proof}\label{proof}

Write

\[
 C=1+\ln q-q,\qquad S=N-w+C,\qquad b=1-\delta.
\]

Then \texttt{alpha=2S/q+b}. Since \texttt{alpha=O(N)}, (1.3) gives

\[
 \ln(n)_\alpha=\alpha N+
 \sum_{j=0}^{\alpha-1}\ln(1-j/n)
 =\alpha N+O(\alpha^2/n),
\]

and hence

\[
 \ln\mu_\alpha
 =\alpha\left(N-\ln\alpha+1-\frac q2(\alpha-1)\right)
  -\frac12\ln(2\pi\alpha)+O(1/N).                       \tag{2.5}
\]

Expanding \texttt{ln\ alpha} at \texttt{2N/q}, uniformly in
\texttt{b\ in\ (0,1{]}}, gives

\[
 N-\ln\alpha+1-\frac q2(\alpha-1)
 =\frac{q\delta}{2}
  +\frac{w-C-bq/2}{N}+O(w^2/N^2).                        \tag{2.6}
\]

Multiplication by \texttt{alpha} proves (2.2). For reference, the
bounded constant is

\[
\begin{split}
 K(\delta)={}&\delta C+\frac q2\delta(1-\delta)-\frac{2C}{q}
 -(1-\delta)\\
 &-\frac12\ln(2\pi)-\frac q2+\frac12\ln q.
\end{split}                                               \tag{2.7}
\]

The exact adjacent-size ratios are

\[
 \frac{\mu_{s+1}}{\mu_s}
 =\frac{n-s}{s+1}2^{-s},\qquad
 \frac{\mu_{s-1}}{\mu_s}
 =\frac{s}{n-s+1}2^{s-1}.                                \tag{2.8}
\]

At \texttt{s=alpha+O(1)}, the first is \texttt{Theta(N/n)} and the
second is \texttt{Theta(n/N)}, uniformly in \texttt{delta}. Applying the
first ratio twice to (2.2) proves (2.3). The minimum of the right side
of (2.2) over the phase is attained, up to a constant factor, at
\texttt{delta=0}; hence
\texttt{mu\_}\allowbreak{}\texttt{alpha\textgreater{}}\allowbreak{}\texttt{=}\allowbreak{}\texttt{c\_}\allowbreak{}\texttt{1N\^{}\{2/}\allowbreak{}\texttt{q-1/}\allowbreak{}\texttt{2\}}.
Applying the second ratio twice proves (2.4). \(\square\)

It follows at once from (2.3) and Markov's inequality that

\[
 \Pr\{\alpha(G_n)>\alpha+1\}=o(1).                       \tag{2.9}
\]

\section{3. Continuous profile roots}\label{continuous-profile-roots}

For a class size \texttt{u}, put

\[
 d_u=2^{\binom u2}u!.
\]

Write \texttt{u=alpha-i}, so that \texttt{i} is the deficit from
\texttt{alpha}. Let

\[
 S_+=\{-1,0,1,2,\ldots\},\qquad S_4=\{2,3,4,5\}.         \tag{3.1}
\]

At finite \texttt{n}, \texttt{S\_+} is cut off at \texttt{i=alpha-1},
since class sizes are positive. For a real class count \texttt{k},
define

\[
 L_{\mathbf k}=nN-n+k-\sum_i k_i\ln(k_id_{\alpha-i}),    \tag{3.2}
\]

and let \texttt{L\_S(n,k)} be its maximum over real
\texttt{k\_i\textgreater{}=0} satisfying

\[
 \sum_{i\in S}k_i=k,\qquad
 \sum_{i\in S}(\alpha-i)k_i=n.                           \tag{3.3}
\]

Set

\[
 T=\alpha-\frac nk.                                      \tag{3.4}
\]

The following lemma supplies all root geometry needed later.

\subsection{Lemma 3.1 (root phase, derivative, and support
comparison)}\label{lemma-3.1-root-phase-derivative-and-support-comparison}

Let

\[
 s_0=\alpha_0-1-\frac2q,\qquad
 T_0=\alpha-s_0=1+\frac2q-\delta.                        \tag{3.5}
\]

For \texttt{S=S\_+} and \texttt{S=S\_4}, the relevant zero of
\texttt{L\_S(n,k)} has

\[
 \frac nk=s_0+O(w/N),\qquad T=T_0+O(w/N).                \tag{3.6}
\]

Uniformly when \texttt{k} is within \texttt{O(n/N\^{}3)} of either zero,

\[
 \frac{\partial L_S}{\partial k}
 =\frac2qN^2+O(Nw).                                      \tag{3.7}
\]

For two supports \texttt{R,S}, at the same feasible \texttt{k},

\[
 \frac{L_R(n,k)-L_S(n,k)}k
 =\mathcal F_{n,R}(T)-\mathcal F_{n,S}(T),               \tag{3.8}
\]

where, uniformly for \texttt{T=T\_0+o(1)},

\[
 \mathcal F_{n,S}(T)\longrightarrow
 \mathcal F_S(T):=\max_{\substack{\sum p_i=1\\\sum ip_i=T}}
 \left[-\sum p_i\ln p_i-\frac q2\sum i^2p_i\right].  \tag{3.9}
\]

The convergence and all error terms are uniform in \texttt{delta}.

\subsubsection{Proof}\label{proof-1}

Put \texttt{s=n/k} and \texttt{p\_i=k\_i/k}. Dividing (3.2) by
\texttt{k} gives

\[
 \frac{L_S(n,n/s)}{n/s}
 =(s-1)N-s+1+\ln s
 +\max_p\left[-\sum p_i\ln p_i-\sum p_i\ln d_{\alpha-i}\right].
                                                               \tag{3.10}
\]

There is an exact affine-plus-curved decomposition

\[
 -\ln d_{\alpha-i}=A_n+B_ni+h_n(i),                      \tag{3.11}
\]

where

\[
 A_n=-\ln d_\alpha,\qquad
 B_n=q\alpha-q/2+\ln\alpha,                              \tag{3.12}
\]

and, for \texttt{i\textgreater{}=0},

\[
 h_n(i)=-\frac q2i^2+\sum_{r=0}^{i-1}\ln(1-r/\alpha).   \tag{3.13}
\]

For \texttt{i=-1}, direct subtraction gives
\texttt{h\_}\allowbreak{}\texttt{n(-1)}\allowbreak{}\texttt{=}\allowbreak{}\texttt{-q/}\allowbreak{}\texttt{2+}\allowbreak{}\texttt{ln(alpha/}\allowbreak{}\texttt{(alpha+}\allowbreak{}\texttt{1)}\allowbreak{}\texttt{)}.
Thus, for fixed \texttt{i},

\[
 h_n(i)=-qi^2/2+O(i^2/\alpha),                           \tag{3.14}
\]

while (3.13) also supplies a uniform Gaussian upper tail. Since the
target mean \texttt{T\_0} stays in the compact interval

\[
 \frac2q\le T_0\le1+\frac2q,                             \tag{3.15}
\]

the Lagrange tilt in the entropy maximum is bounded. Its optimizer is
proportional to
\texttt{exp(lambda\ i+}\allowbreak{}\texttt{h\_}\allowbreak{}\texttt{n(i)}\allowbreak{}\texttt{)},
and (3.13) makes the omitted tail uniformly summable. Dominated
convergence proves (3.9). The affine terms in (3.11) prove the exact
difference identity (3.8).

Substitute (3.11) into (3.10). If \texttt{T=alpha-s}, the scalar part at
\texttt{s=s\_0} is, by Stirling,

\[
 s_0\left(N-\ln\alpha-\frac q2s_0+\frac q2\right)-N
 +T_0+1+\ln s_0+\frac q2T_0^2
 -\frac12\ln(2\pi\alpha)+O(1/N).                        \tag{3.15a}
\]

Now
\texttt{q\ s\_}\allowbreak{}\texttt{0/}\allowbreak{}\texttt{2=}\allowbreak{}\texttt{N-w+}\allowbreak{}\texttt{ln\ q-q}
exactly, while
\texttt{ln\ alpha=}\allowbreak{}\texttt{ln\ s\_}\allowbreak{}\texttt{0+}\allowbreak{}\texttt{T\_}\allowbreak{}\texttt{0/}\allowbreak{}\texttt{s\_}\allowbreak{}\texttt{0+}\allowbreak{}\texttt{O(1/}\allowbreak{}\texttt{N\^{}2)}
and
\texttt{ln\ s\_}\allowbreak{}\texttt{0=}\allowbreak{}\texttt{w+}\allowbreak{}\texttt{q-ln\ q+}\allowbreak{}\texttt{O(w/}\allowbreak{}\texttt{N)}.
These relations reduce (3.15a) to \texttt{O(w)}, uniformly in the phase.
The entropy term is \texttt{O(1)} by the Gaussian tail just proved.
Hence

\[
 L_S(n,n/s_0)/(n/s_0)=O(w).                              \tag{3.16}
\]

Differentiating (3.10) by the envelope theorem, and using

\[
 B_n=2N-w+O(1),                                          \tag{3.17}
\]

gives

\[
 \frac{d}{ds}\left\{\frac{L_S(n,n/s)}{n/s}\right\}
 =-N+O(w)                                                \tag{3.18}
\]

in a fixed neighbourhood of \texttt{s\_0}. Equations (3.16)--(3.18) give
the unique local zero and (3.6). Finally, if the expression in braces is
\texttt{Psi(s)}, then \texttt{L=k\ Psi(s)} and \texttt{ds/dk=-s/k}; at a
zero,

\[
 \frac{dL}{dk}=-s\Psi'(s)=\frac2qN^2+O(Nw).              \tag{3.19}
\]

Changing \texttt{k} by \texttt{O(n/N\^{}3)} changes \texttt{s} by only
\texttt{O(1/N)}, so the same estimate is uniform throughout the stated
window. \(\square\)

\section{\texorpdfstring{4. A valid unrestricted lower location for
\texttt{chi}}{4. A valid unrestricted lower location for chi}}\label{a-valid-unrestricted-lower-location-for-chi}

Let \texttt{r\_+(n)} be the zero in Lemma 3.1 for \texttt{S\_+}, and put

\[
 k_\chi^-:=\lfloor r_+(n)\rfloor-\lceil N\rceil.         \tag{4.1}
\]

For an integer profile, direct enumeration gives

\[
 \mathbb E X_{\mathbf k}
 =\frac{n!}{\prod_i((\alpha-i)!)^{k_i}k_i!}
  2^{-\sum_i k_i\binom{\alpha-i}{2}}.                   \tag{4.2}
\]

There are at most
\texttt{(n+}\allowbreak{}\texttt{1)}\allowbreak{}\texttt{\^{}(alpha+}\allowbreak{}\texttt{1)}\allowbreak{}\texttt{=}\allowbreak{}\texttt{exp(O(N\^{}2)}\allowbreak{}\texttt{)}
bounded profiles. Applying (1.3) to (4.2), uniformly even when some
\texttt{k\_i=0}, yields

\[
 \ln E_{n,k,\alpha+1}\le L_+(n,k)+O(N^2),                \tag{4.3}
\]

where the left side is the logarithm of the expected total number of
unordered, exactly \texttt{k}-part, \texttt{(alpha+1)}-bounded
colourings. By Lemma 3.1,

\[
 L_+(n,k_\chi^-)\le-cN^3,                                \tag{4.4}
\]

so the expectation in (4.3) is \texttt{o(1)}. On the event (2.9), any
colouring with at most \texttt{k\_chi\^{}-} nonempty parts can be
refined, by splitting parts, to exactly \texttt{k\_chi\^{}-} nonempty
independent parts, all of size at most \texttt{alpha+1}. Therefore

\[
 \Pr\{\chi(G_n)\le k_\chi^-\}=o(1).                     \tag{4.5}
\]

This is an unrestricted chromatic lower bound; the size cap has been
removed probabilistically by (2.9).

\section{5. The four-size signed first-moment
advantage}\label{the-four-size-signed-first-moment-advantage}

The optimizer in (3.9) has the form

\[
 p_i=\frac{e^{\lambda i-qi^2/2}}
           {\sum_{j\in S}e^{\lambda j-qj^2/2}},          \tag{5.1}
\]

where \texttt{lambda} is chosen to give mean \texttt{T}. Define

\[
 D_4(\delta)=\mathcal F_{S_+}(T_0)
              -\mathcal F_{S_4}(T_0),\qquad
 \gamma_4=\ln\frac{200}{153}.                            \tag{5.2}
\]

\subsection{Lemma 5.1 (uniform entropy
certificate)}\label{lemma-5.1-uniform-entropy-certificate}

For every \texttt{0\textless{}=delta\textless{}=1},

\[
 0\le D_4(\delta)<\ln\frac{153}{100},\qquad
 q-D_4(\delta)>\gamma_4.                                 \tag{5.3}
\]

\subsubsection{Proof}\label{proof-2}

Let \texttt{lambda\_4} be the tilt for \texttt{S\_4}. At
\texttt{lambda=2q}, shift the index by two. The four-size mean is no
larger than

\[
 2+\frac{\sum_{j\ge0}j2^{-j^2/2}}
          {\sum_{j\ge0}2^{-j^2/2}}<2+\frac45<\frac2q.    \tag{5.4}
\]

For the first strict inequality, retain the first three terms of the
denominator, sum the numerator through \texttt{j=3}, and bound its
remaining geometric tail by \texttt{1/60}; the elementary bounds
\texttt{2/3\textless{}q\textless{}5/7} prove the second. At
\texttt{lambda=9q/2}, put \texttt{t=2\^{}\{-1/8\}}. The four weights are
proportional to

\[
 t^{25},\quad t^9,\quad t,\quad t.
\]

Their mean is greater than four because

\[
 t-t^9-2t^{25}=t(1-t^8-2t^{24})=t/4>0,                 \tag{5.5}
\]

whereas \texttt{1+2/q\textless{}4}. Hence

\[
 2q<\lambda_4<9q/2.                                     \tag{5.6}
\]

For a tilt \texttt{lambda}, let \texttt{L(lambda)} be the total weight
of deficits \texttt{-1,0,1} divided by the weight of deficits
\texttt{2,3,4,5}, and let \texttt{H(lambda)} be the analogous ratio for
deficits at least six. The first ratio decreases and the second
increases with \texttt{lambda}. Directly summing the displayed finite
terms and bounding the Gaussian tails gives

\[
 L(2q)<\frac{51}{100},\quad H(3q)<\frac1{50},\quad
 L(3q)<\frac3{25},\quad H(9q/2)<\frac14.                 \tag{5.7}
\]

Here are explicit checks. At \texttt{2q}, with \texttt{b=2\^{}\{-1/2\}},

\[
 L(2q)=\frac{b+1/4+b/16}{1+b+1/4+b/16}<51/100.
\]

At \texttt{3q}, after centring at deficit three, the denominator is
\texttt{5/4+2b}; the high numerator is at most
\texttt{b/16+1/256+1/3968}, and the low numerator is
\texttt{1/256+b/16+1/4}. The inequalities follow from
\texttt{7/10\textless{}b\textless{}71/100}. Finally, after division by
\texttt{t}, the tail above deficit five at \texttt{9q/2} begins with
\texttt{t\^{}8,t\^{}24,t\^{}48}; the rest is below \texttt{1/60}, and

\[
 3(t^8+t^{24})+4/60<2,                                  \tag{5.8}
\]

which is equivalent to \texttt{H(9q/2)\textless{}1/4}.

If \texttt{lambda\_4\textless{}=3q}, monotonicity and (5.7) give
\texttt{L+H\textless{}53/100}; if \texttt{lambda\_4\textgreater{}=3q},
they give \texttt{L+H\textless{}37/100}. Evaluating the dual function
for \texttt{S\_+} at the \texttt{S\_4} tilt yields

\[
 D_4\le\ln(1+L(\lambda_4)+H(\lambda_4))<\ln(153/100).
\]

Since \texttt{S\_4} is a subset of \texttt{S\_+}, the loss is
nonnegative. Subtracting from \texttt{q=ln2} completes the proof.
\(\square\)

For a fixed partition into \texttt{k} classes, assigning each class the
declaration ``independent'' or ``complete'' multiplies its first moment
by \texttt{2\^{}k}: every declaration prescribes all internal edge bits,
and all \texttt{2\^{}k} declarations have probability
\texttt{2\^{}\{-sum\ binom(u,}\allowbreak{}\texttt{2)}\allowbreak{}\texttt{\}}.
Let \texttt{r\_4\^{}\{co\}} be the zero of

\[
 L_{S_4}(n,k)+qk.                                       \tag{5.9}
\]

At \texttt{k=r\_+}, Lemma 3.1 and (5.2) give

\[
 L_{S_4}(n,r_+)+qr_+
 =r_+\{q-D_4(\delta)+o(1)\}.                             \tag{5.10}
\]

Since \texttt{r\_+=(q/2+o(1))n/N}, the mean-value theorem and (3.7) give

\[
 r_+-r_4^{co}
 =\left(\frac{q^2}{4}\{q-D_4(\delta)\}+o(1)\right)
   \frac n{N^3}.                                        \tag{5.11}
\]

In particular, for all sufficiently large \texttt{n},

\[
 r_+-r_4^{co}\ge\frac{q^2\gamma_4}{8}\frac n{N^3}.     \tag{5.12}
\]

Choose the midpoint integer

\[
 k_{co}=\left\lceil\frac{r_4^{co}+r_+}{2}\right\rceil. \tag{5.13}
\]

At this exact \texttt{k}, let \texttt{p\_i\^{}\{(n)\}} be the exact
finite-\texttt{n} maximizer of \texttt{L\_\{S\_4\}}. Its Lagrange
equations give

\[
 p_i^{(n)}\ \propto\ d_{\alpha-i}^{-1}e^{\lambda_ni},  \tag{5.14}
\]

subject to mean \texttt{alpha-n/k\_co}. Lemma 3.1 makes these
proportions converge uniformly to (5.1). Equation (5.6) also shows
directly that the four limiting weights differ by at most a fixed
factor. Thus

\[
 \min_{2\le i\le5}p_i^{(n)}\ge c_p>0                    \tag{5.15}
\]

uniformly in the phase.

Round \texttt{k\_co\ p\_i\^{}(n)} to integers. If the errors in the
total-count and total-deficit constraints are \texttt{e\_0,e\_1=O(1)},
add

\[
 \Delta k_2=e_1-3e_0,\qquad
 \Delta k_3=2e_0-e_1.                                   \tag{5.16}
\]

This enforces both constraints exactly, changes only \texttt{O(1)}
counts, and preserves positivity by (5.15). Because the displacement is
tangent to both constraints and starts at the exact finite optimizer,
its first variation vanishes and its Hessian cost is
\texttt{O(1/k\_co)}. Applying (1.3) to the five factorials gives an
exact integer vector

\[
 \mathbf k=(k_2,k_3,k_4,k_5),\qquad
 u_i=\alpha-i,\qquad k_i=\Theta(n/N),                    \tag{5.17}
\]

whose signed first moment is

\[
 \mathbb EZ_{\mathbf k}^{sgn}
 =2^{k_{co}}\frac{n!}{\prod_i(u_i!)^{k_i}k_i!}
   2^{-\sum_i k_i\binom{u_i}{2}},                        \tag{5.18}
\]

and

\[
 \ln\mathbb EZ_{\mathbf k}^{sgn}
 =L_{S_4}(n,k_{co})+qk_{co}+O(N)
 \ge c_Zk_{co}                                          \tag{5.19}
\]

for a phase-independent \texttt{c\_Z\textgreater{}0}. The last
inequality follows from the midpoint distance, (3.7), and (5.12). Every
\texttt{u\_i} tends to infinity, so the independent and complete
declarations of one part are disjoint. Thus
\texttt{Z\_}\allowbreak{}\texttt{\{\textbackslash{}mathbf\ k\}\^{}\{sgn\}\textgreater{}}\allowbreak{}\texttt{0}
implies an actual cocolouring with \texttt{k\_co} parts.

Finally, (4.1), (5.12), and (5.13) give

\[
 k_\chi^- -k_{co}
 \ge\left(\frac{q^2\gamma_4}{16}-o(1)\right)
       \frac n{N^3}.                                    \tag{5.20}
\]

\section{6. Exact signed second-moment
representation}\label{exact-signed-second-moment-representation}

Temporarily label all slots within each of the four types. This
multiplies the unordered witness by the deterministic factor
\texttt{prod\_i\ k\_i!} and hence does not change its normalized second
moment. Choose two independent uniform ordered partitions of the
profile. If their slots are indexed by \texttt{a} and \texttt{b}, put

\[
 r_{ab}=|V_a\cap W_b|.                                   \tag{6.1}
\]

The overlap matrix has the exact law

\[
 p(r)=\frac{\prod_a s_a!\prod_b t_b!}
            {n!\prod_{a,b}r_{ab}!},                      \tag{6.2}
\]

where here both degree lists are the multiset given by (5.17).
Equivalently, place \texttt{s\_a} stubs at each row slot and
\texttt{t\_b} stubs at each column slot and take a uniform perfect
matching of the two \texttt{n}-stub sets.

Let \texttt{H(r)} be the simple bipartite graph whose edges are cells
with \texttt{r\_ab\textgreater{}=2}, omitting isolated slot vertices.
Write

\[
 W=\sum_{a,b}\binom{r_{ab}}2,\qquad
 \beta(H)=|E(H)|-|V(H)|+c(H).                             \tag{6.3}
\]

\subsection{Lemma 6.1 (exact sign sum)}\label{lemma-6.1-exact-sign-sum}

The normalized local factor of the pair of partitions is

\[
 A_\zeta(r)=2^{W+c(H)-|V(H)|}
 =\left(\prod_{a,b}g(r_{ab})\right)2^{\beta(H)},          \tag{6.4}
\]

where

\[
 g(0)=g(1)=g(2)=1,\qquad
 g(x)=2^{\binom x2-1}\quad(x\ge3).                       \tag{6.5}
\]

Consequently

\[
 \frac{\mathbb E(Z_{\mathbf k}^{sgn})^2}
      {(\mathbb EZ_{\mathbf k}^{sgn})^2}
 =\sum_rp(r)A_\zeta(r).                                  \tag{6.6}
\]

\subsubsection{Proof}\label{proof-3}

Let
\texttt{B=}\allowbreak{}\texttt{sum\_}\allowbreak{}\texttt{a\ binom(s\_}\allowbreak{}\texttt{a,}\allowbreak{}\texttt{2)},
which is the number of internal edge bits prescribed by one partition. A
row sign and a column sign are compatible on a cell of size at least two
exactly when they agree. Therefore signs must be constant on every
component of \texttt{H}. There are

\[
 2^{2k_{co}-|V(H)|+c(H)}
\]

compatible pairs of sign assignments, including the free signs on slots
isolated from \texttt{H}. For each compatible pair, \texttt{W} edge bits
are prescribed twice, so the probability of all constraints is
\texttt{2\^{}\{-(2B-W)\}}. Division by the square of the one-partition
signed probability
\texttt{(2\^{}\{k\_}\allowbreak{}\texttt{co\}2\^{}\{-B\})}\allowbreak{}\texttt{\^{}2}
gives the first expression in (6.4).

For every support edge, split \texttt{binom(r,2)} as
\texttt{(binom(r,2)-1)+1}. Since
\texttt{\textbar{}E\textbar{}-\textbar{}V\textbar{}+c=beta}, this gives
the second expression. Averaging over (6.2) proves (6.6). \(\square\)

The identity

\[
 2^{\beta(H)}
 =\#\{F\subseteq E(H):\deg_F(v)\text{ is even for every }v\} \tag{6.7}
\]

will be used repeatedly. It is the elementary fact that the even
subgraphs form the binary cycle space of dimension \texttt{beta(H)}.

We also need one exact configuration-model estimate.

\subsection{Lemma 6.2 (joint prescribed-cell
bound)}\label{lemma-6.2-joint-prescribed-cell-bound}

In a bipartite configuration model of total degree \texttt{m\_0}, with
row degrees \texttt{d\_a} and column degrees
\texttt{d\textquotesingle{}\_b}, prescribe distinct cells \texttt{D} and
demands \texttt{x\_ab\textgreater{}=1}. Put

\[
 x=\sum_Dx_{ab},\quad D_a=\sum_bx_{ab},\quad
 D'_b=\sum_ax_{ab}.
\]

Then

\[
 \Pr\{r_{ab}\ge x_{ab}\ \forall ab\in D\}
 \le\frac{\prod_a(d_a)_{D_a}\prod_b(d'_b)_{D'_b}}
          {(m_0)_x\prod_{ab\in D}x_{ab}!}.               \tag{6.8}
\]

In particular, with
\texttt{theta\_}\allowbreak{}\texttt{ab=}\allowbreak{}\texttt{e\ d\_}\allowbreak{}\texttt{a\ d\textquotesingle{}\_}\allowbreak{}\texttt{b/}\allowbreak{}\texttt{m\_}\allowbreak{}\texttt{0},
the right side is at most

\[
 \prod_{ab\in D}\frac{\theta_{ab}^{x_{ab}}}{x_{ab}!}.  \tag{6.9}
\]

\subsubsection{Proof}\label{proof-4}

Choose the required row and column stubs and biject them inside each
prescribed cell. This gives the numerator of (6.8); any fixed set of
\texttt{x} paired stubs occurs with probability \texttt{1/(m\_0)\_x}. A
union bound proves (6.8). Next use \texttt{(d)\_r\textless{}=d\^{}r} and

\[
 (m_0)_x=x!\binom{m_0}{x}\ge(m_0/e)^x,                  \tag{6.10}
\]

which proves (6.9). Notice that (6.8) retains the single global falling
factorial before (6.9) is taken. \(\square\)

\section{7. Exact partial diagonals}\label{exact-partial-diagonals}

For a selected common subprofile \texttt{ell=(ell\_i)}, put

\[
 L=\sum_i\ell_i,\qquad
 m=\sum_i u_i\ell_i,\qquad
 M=\sum_i\binom{u_i}{2}\ell_i.                           \tag{7.1}
\]

The signed first moment of such partial partitions occupying any
\texttt{m} vertices of \texttt{{[}n{]}} is

\[
 Y_{\ell}^{sgn}
 =2^L\frac{n!}{(n-m)!\prod_i(u_i!)^{\ell_i}\ell_i!}
       2^{-M}.                                           \tag{7.2}
\]

Summing the possible matchings between selected row and column slots
gives the exact common-subprofile contribution

\[
 D(\ell)=\frac{\prod_i\binom{k_i}{\ell_i}^2}
                 {Y_\ell^{sgn}}.                         \tag{7.3}
\]

In particular, \texttt{D(0)=1} and \texttt{D(k)=1/E\ Z\_k\^{}sgn}.
Adding one common block gives

\[
 \frac{D(\ell+e_i)}{D(\ell)}
 =\frac{(k_i-\ell_i)^2}
        {2(\ell_i+1)\mu_{u_i}(n-m)}.                     \tag{7.4}
\]

At the opposite endpoint, put \texttt{h=k-ell},
\texttt{v=sum\_i\ u\_i\ h\_i}, and let \texttt{Z\_h\^{}sgn(v)} be the
signed first moment of a complete profile \texttt{h} on \texttt{v}
vertices. Exact factorial cancellation gives

\[
 D(k-h)=\frac{B(h)}{\mathbb EZ_k^{sgn}},\qquad
 B(h)=\left(\prod_i\binom{k_i}{h_i}\right)Z_h^{sgn}(v), \tag{7.5}
\]

and

\[
 \frac{B(h+e_i)}{B(h)}
 =\frac{2(k_i-h_i)\mu_{u_i}(v+u_i)}{(h_i+1)^2}.          \tag{7.6}
\]

\subsection{Lemma 7.1 (all common
subprofiles)}\label{lemma-7.1-all-common-subprofiles}

For the exact midpoint profile,

\[
 \sum_{0\le\ell_i\le k_i}D(\ell)=1+o(1).               \tag{7.7}
\]

\subsubsection{Proof: the empty corner}\label{proof-the-empty-corner}

Let

\[
 \eta=\frac{w}{32N},\qquad
 \lambda_i=\frac{k_i^2}{2\mu_{u_i}(n)},\qquad
 \Lambda=\sum_i\lambda_i.                               \tag{7.8}
\]

Equations (2.4), (2.8), and \texttt{k\_i=Theta(n/N)} imply

\[
 \Lambda=O(N^{-(2/q-1/2)})=o(1).                        \tag{7.9}
\]

If the final selected mass is at most \texttt{eta\ n}, every
intermediate mass \texttt{m\textquotesingle{}} in (7.4) satisfies

\[
 \frac{\mu_{u_i}(n)}{\mu_{u_i}(n-m')}
 \le(1-2\eta)^{-u_i}\le N^{3/8}.                        \tag{7.10}
\]

Iterating (7.4), in any order, yields

\[
 D(\ell)\le\prod_i\frac{(N^{3/8}\lambda_i)^{\ell_i}}
                              {\ell_i!}.
\]

Thus

\[
 1\le\sum_{m\le\eta n}D(\ell)
 \le\exp(N^{3/8}\Lambda)=1+o(1).                       \tag{7.11}
\]

\subsubsection{Proof: the central range}\label{proof-the-central-range}

Write

\[\begin{gathered}
p_i=k_i/k_{co},\quad y_i=\ell_i/k_{co},\quad
 Y=\sum_i y_i,\quad I=\sum_i i y_i,\quad\\
z_i=p_i-y_i,\quad R=\sum_i z_i=1-Y,\quad
 I_r=\sum_i i z_i=T-I.
\end{gathered}
\tag{7.12}\]

The residual vertex fraction is

\[
 \rho=\frac{n-m}{n}
 =R+\frac{I-TY}{\alpha-T}.                              \tag{7.13}
\]

Uniform Stirling estimates in (7.3), retaining zeros, give

\[
\begin{split}
 \ln D(\ell)={}&n\rho\ln\rho\\
 &+k_{co}\sum_i\{2p_i\ln p_i-2(p_i-y_i)\ln(p_i-y_i)
                  -y_i\ln y_i-y_i+y_iE_i\}+O(N),        \end{split}
\tag{7.14}
\]

where

\[
 E_i=\ln k_{co}+\ln(u_i!)+u_i-u_iN
       +q\binom{u_i}{2}-q.                               \tag{7.15}
\]

Exact subtraction and the definition of \texttt{alpha\_0} give

\[
 E_{i+1}-E_i=N-\ln(\alpha-i)-q(\alpha-i)+q-1
            =-q\alpha/2+O(1).                            \tag{7.16}
\]

Let \texttt{bar\ E=sum\ p\_iE\_i}. Applying Stirling to the complete
moment (5.18) gives

\[
 -\frac1{k_{co}}\ln\mathbb EZ_k^{sgn}
 =\sum_i p_i\ln p_i-1+\bar E+o(1).                      \tag{7.17}
\]

The positive margin (5.19) bounds \texttt{bar\ E} above. Since
\texttt{sum\ p\_i(i-T)=0}, (7.16) implies

\[
 \sum_i y_iE_i\le\frac{q\alpha}{2}(TY-I)+O(Y).          \tag{7.18}
\]

The remaining entropy expression in (7.14) is \texttt{O(Y\ ln(e/Y))},
uniformly because every \texttt{p\_i} is bounded below. If
\texttt{rho\textgreater{}=1/32}, then \texttt{R\textgreater{}=1/64} for
large \texttt{n}, and (7.13) gives

\[
 n\rho\ln\rho=k_{co}\alpha R\ln R+O(k_{co}Y).           \tag{7.19}
\]

Consequently

\[
 \ln D(\ell)
 \le k_{co}\alpha\Phi_T(z)
     +Ck_{co}Y\ln(e/Y)+CN,                               \tag{7.20}
\]

where

\[
 \Phi_T(z)=R\ln R+\frac q2(I_r-TR).                     \tag{7.21}
\]

This rate is uniformly negative away from its two corners. Indeed,

\[
 I_r-TR\le(5-T)R,\qquad
 I_r-TR=\sum_i(T-i)y_i\le(T-2)(1-R),                    \tag{7.22}
\]

and hence

\[
 \Phi_T\le R\{\ln R+q(5-T)/2\},
 \quad
 \Phi_T\le R\ln R+q(T-2)(1-R)/2.                       \tag{7.23}
\]

On \texttt{R\textless{}=47/100}, the first bound is at most
\texttt{-Y/5000}, using \texttt{q\textless{}0.6932}; on
\texttt{R\textgreater{}=47/100}, convexity of

\[
 R\ln R+(1-q/2+1/200)(1-R)
\]

between \texttt{47/100} and \texttt{1} gives the same conclusion, in
fact with \texttt{1/200} in place of \texttt{1/5000}. Thus, whenever

\[
 m>\eta n,\qquad n-m>n/32,                               \tag{7.24}
\]

we have \texttt{Y\textgreater{}=eta/2}, and the leading negative term in
(7.20) dominates the entropy error. Uniformly in this range,

\[
 D(\ell)\le e^{-c k_{co}w}.                              \tag{7.25}
\]

There are only \texttt{(k\_co+1)\^{}4} vectors, so their total is
\texttt{o(1)}.

\subsubsection{Proof: the full corner}\label{proof-the-full-corner}

It remains to take \texttt{v=n-m\textless{}=n/32}. Uniformly for the
four sizes,
\texttt{mu\_}\allowbreak{}\texttt{\{u\_}\allowbreak{}\texttt{i\}(n)}\allowbreak{}\texttt{\textless{}}\allowbreak{}\texttt{=}\allowbreak{}\texttt{n\^{}\{6+}\allowbreak{}\texttt{o(1)}\allowbreak{}\texttt{\}}
by (2.2) and (2.8), while

\[
 \mu_{u_i}(v)\le\mu_{u_i}(n)(v/n)^{u_i}\le n^{-4+o(1)}. \tag{7.26}
\]

Here \texttt{u\_i=(2/q+o(1))N} and
\texttt{32\^{}\{-u\_}\allowbreak{}\texttt{i\}=}\allowbreak{}\texttt{n\^{}\{-10+}\allowbreak{}\texttt{o(1)}\allowbreak{}\texttt{\}}.
In any order of adding residual blocks, (7.6) is therefore below one for
all sufficiently large \texttt{n}; hence \texttt{B(h)\textless{}=1}.
Equations (5.19) and (7.5) give

\[
 D(k-h)\le e^{-c_Zk_{co}}.                               \tag{7.27}
\]

Summing \texttt{(k\_co+1)\^{}4} residual vectors proves an \texttt{o(1)}
contribution. Together with (7.11) and (7.25), this proves (7.7).
\(\square\)

\section{8. Canonical high cells and dense endpoint
transportation}\label{canonical-high-cells-and-dense-endpoint-transportation}

Put

\[
 U=\alpha-2,\qquad R_0=\lfloor U/2\rfloor.                \tag{8.1}
\]

In every overlap matrix, the cells of multiplicity greater than
\texttt{R\_0} form a matching: two such cells in one row or column would
use more than \texttt{U} stubs. Let this canonical matching be

\[
 M=\{a_rb_r:1\le r\le h\},\qquad
 r_{a_rb_r}=j_r>R_0,\qquad J=\sum_rj_r.                  \tag{8.2}
\]

Selecting the exact stub pairs in these cells has incidence

\[
 \pi(M,j)=
 \frac{\prod_r(s_{a_r})_{j_r}(t_{b_r})_{j_r}}
      {(n)_J\prod_rj_r!}.                                \tag{8.3}
\]

Conditional on the exposed pairs, the remaining matching is a uniform
configuration model with total degree \texttt{n-J} and the residual row
and column degrees. It is restricted to put no further pair in a cell of
\texttt{M} and no cell above \texttt{R\_0}. Multiplication of (8.3) by
the exact residual contingency-table law cancels to (6.2). Thus every
overlap matrix occurs exactly once.

We first sum the bare high matching, postponing all local and
topological factors in the capped residual model. Index the four sizes
as

\[
 u_i=a-i,\qquad 0\le i\le3,\qquad a=\alpha-2.             \tag{8.4}
\]

For a typed full-containment matrix \texttt{L=(ell\_ij)}, put

\[
 r_i=\sum_j\ell_{ij},\quad c_j=\sum_i\ell_{ij},\quad
 x_{ij}=\min(u_i,u_j),\quad J(L)=\sum_{ij}x_{ij}\ell_{ij}. \tag{8.5}
\]

Its exact incidence, including its local signed rewards but not residual
attachments, is

\[
 W(L)=
 \frac{\prod_i(k_i)_{r_i}\prod_j(k_j)_{c_j}}
      {\prod_{ij}\ell_{ij}!}\frac1{(n)_{J(L)}}
 \prod_{ij}\left[
  \frac{(u_i)_{x_{ij}}(u_j)_{x_{ij}}}{x_{ij}!}g(x_{ij})
 \right]^{\ell_{ij}}.                                   \tag{8.6}
\]

For a margin vector \texttt{r}, define

\[
 D(r)=
 \frac{\prod_i(k_i)_{r_i}^2}{\prod_i r_i!}
 \frac{\prod_i[u_i!g(u_i)]^{r_i}}{(n)_{m_r}},\qquad
 m_r=\sum_i u_ir_i.                                     \tag{8.7}
\]

This is exactly the common-subprofile weight (7.3).

\subsection{Lemma 8.1 (geometric-mean transportation
comparison)}\label{lemma-8.1-geometric-mean-transportation-comparison}

For every feasible \texttt{L},

\[
 W(L)\le\sqrt{D(r)D(c)}
 \frac{\sqrt{\prod_i r_i!c_i!}}{\prod_{ij}\ell_{ij}!}
 \prod_{ij}Q_{ij}^{\ell_{ij}},                          \tag{8.8}
\]

where \texttt{Q\_ii=1} and, if the two sizes are \texttt{t=s+d},
\texttt{1\textless{}=d\textless{}=3},

\[
 Q_{ij}=(n+1)^{d/2}\frac{\sqrt{(t)_d}}{d!}
       2^{-\{ds+\binom d2\}/2}
 \le\frac{\eta_n^d}{d!},                               \tag{8.9}
\]

with

\[
 \eta_n=2^{3/2}\sqrt{\frac{(n+1)a}{2^a}}
       =O(N^{3/2}/\sqrt n)=o(1).                         \tag{8.10}
\]

\subsubsection{Proof}\label{proof-5}

For integer \texttt{x}, let \texttt{f(x)=ln(n)\_x} and interpolate
linearly between successive integers. Its successive slopes are
\texttt{ln(n-x)}, which decrease; therefore \texttt{f} is concave, and
every slope is at most \texttt{ln(n+1)}. Moreover,

\[
 \frac{m_r+m_c}{2}
 =J(L)+\frac12\sum_{ij}|i-j|\ell_{ij}.                  \tag{8.11}
\]

Concavity and \texttt{f\textquotesingle{}(x)\textless{}=ln(n+1)} give

\[
 \frac{\sqrt{(n)_{m_r}(n)_{m_c}}}{(n)_{J(L)}}
 \le(n+1)^{\frac12\sum_{ij}|i-j|\ell_{ij}}.            \tag{8.12}
\]

For \texttt{t=s+d}, put \texttt{b\_u=u!g(u)}. A direct calculation gives

\[
 \frac{(t)_s(s)_s}{s!}g(s)=b_s\binom td,\qquad
 \frac{b_t}{b_s}=(t)_d2^{ds+\binom d2}.                 \tag{8.13}
\]

Thus the local factor divided by \texttt{sqrt(b\_sb\_t)} is

\[
 \frac{\sqrt{(t)_d}}{d!}2^{-\{ds+\binom d2\}/2}.        \tag{8.14}
\]

Dividing (8.6) by \texttt{sqrt(D(r)D(c))}, then using (8.12) and (8.14),
proves (8.8)--(8.9). Finally \texttt{t\textless{}=a},
\texttt{s\textgreater{}=a-3}, and

\[
 2^a=\Theta(n^2/N^2)                                    \tag{8.15}
\]

uniformly in the phase. This proves (8.10). \(\square\)

\subsection{Lemma 8.2 (sum of all endpoint
tables)}\label{lemma-8.2-sum-of-all-endpoint-tables}

\[
 \sum_LW(L)\le
 \exp\{O(\sqrt{nN})+O(N)\}\sum_rD(r)
 =\exp\{O(\sqrt{nN})+O(N)\}.                            \tag{8.16}
\]

\subsubsection{Proof}\label{proof-6}

For fixed margins, write

\[
 A_L=\frac{\prod_i r_i!}{\prod_{ij}\ell_{ij}!},\qquad
 C_L=\frac{\prod_j c_j!}{\prod_{ij}\ell_{ij}!}.        \tag{8.17}
\]

Cauchy's inequality gives

\[
 \sum_L\sqrt{A_LC_L}\,Q^L
 \le\left(\sum_LA_LQ^L\right)^{1/2}
      \left(\sum_LC_LQ^L\right)^{1/2}.                 \tag{8.18}
\]

Dropping the column constraints in the first sum and the row constraints
in the second produces two exact multinomial expansions:

\[
 \sum_LA_LQ^L\le\prod_i\left(\sum_jQ_{ij}\right)^{r_i},
 \quad
 \sum_LC_LQ^L\le\prod_j\left(\sum_iQ_{ij}\right)^{c_j}.
                                                               \tag{8.19}
\]

Every row and column sum of \texttt{Q} is at most \texttt{1+C\ eta\_n},
and the total margin is at most \texttt{k\_co}. Hence the fixed-margin
cost is
\texttt{exp(O(eta\_}\allowbreak{}\texttt{n\ k\_}\allowbreak{}\texttt{co)}\allowbreak{}\texttt{)}.
There are only \texttt{O(k\_co\^{}4)} margin vectors, and

\[
 \left(\sum_r\sqrt{D(r)}\right)^2
 \le O(k_{co}^4)\sum_rD(r).                              \tag{8.20}
\]

Use Lemma 7.1 and
\texttt{eta\_}\allowbreak{}\texttt{n\ k\_}\allowbreak{}\texttt{co=}\allowbreak{}\texttt{O(sqrt(nN)}\allowbreak{}\texttt{)}
to finish. \(\square\)

\subsection{Lemma 8.3 (all nonendpoint high
multiplicities)}\label{lemma-8.3-all-nonendpoint-high-multiplicities}

The full-containment sum (8.16) remains \texttt{exp(o(n/N\^{}4))} after
all high-cell multiplicities
\texttt{R\_}\allowbreak{}\texttt{0\textless{}}\allowbreak{}\texttt{j\textless{}}\allowbreak{}\texttt{min(u\_}\allowbreak{}\texttt{i,}\allowbreak{}\texttt{u\_}\allowbreak{}\texttt{j)}
are inserted.

\subsubsection{Proof}\label{proof-7}

Let the smaller slot size be \texttt{m}, the larger \texttt{m+d}, with
\texttt{0\textless{}=d\textless{}=3}. Replacing the endpoint
multiplicity \texttt{m} by \texttt{m-e} has exact local ratio

\[
 R_{m,d}(e)=
 \frac{\binom me}{(d+1)\cdots(d+e)}
 2^{-em+e(e+1)/2}.                                      \tag{8.21}
\]

For many decorated cells, with total deficit \texttt{Q}, the ratio of
the one global denominators is

\[
 \frac{(n)_{J_0}}{(n)_{J_0-Q}}
 =(n-J_0+Q)_Q\le n^Q.                                   \tag{8.22}
\]

For \texttt{1\textless{}=e\textless{}m/4}, the consecutive ratio of the
summands \texttt{n\^{}eR\_\{m,d\}(e)} is

\[
 \frac{n(m-e)2^{-m+e+1}}{(e+1)(d+e+1)}.                 \tag{8.23}
\]

The ratio of (8.23) at \texttt{e+1} to its value at \texttt{e} is

\[
 2\frac{m-e-1}{m-e}\frac{e+1}{e+2}
  \frac{d+e+1}{d+e+2},                                  \tag{8.24}
\]

which increases for \texttt{e\textless{}=m/4} and
\texttt{d\textless{}=3}. Thus (8.23) is largest at an endpoint. At
\texttt{e=0} it is \texttt{O(N\^{}3/n)}, by (8.15), and at
\texttt{e=floor(m/4)} it is \texttt{n\^{}\{-1/2+o(1)\}}. The series
therefore decreases geometrically and

\[
 \sum_{1\le e<m/4}n^eR_{m,d}(e)=O(N^3/n).                \tag{8.25}
\]

The multinomial choices among identical typed cells turn this into a
factor

\[
 (1+O(N^3/n))^{k_{co}}=\exp(O(N^2)).                     \tag{8.26}
\]

It remains to cover the middle strip, rather than silently importing an
equitable-profile estimate. First expose a near-containment skeleton and
let \texttt{m\_0} be its residual stub mass. If
\texttt{m\_0\textgreater{}=n/N\^{}6}, Lemma 6.2, before dropping the
matching constraints, bounds the entire remaining middle-cell expansion
by \texttt{exp(Xi\_4)}, where

\[
 \Xi_4\le k_{co}^2
 \sum_{a/2<j\le3a/4+O(1)}
  g(j)\frac{(ea^2/m_0)^j}{j!}.                           \tag{8.27}
\]

Put \texttt{L=log\_2\ n} and \texttt{j=xL}. Uniformly for
\texttt{1+}\allowbreak{}\texttt{o(1)}\allowbreak{}\texttt{\textless{}}\allowbreak{}\texttt{=}\allowbreak{}\texttt{x\textless{}}\allowbreak{}\texttt{=}\allowbreak{}\texttt{3/}\allowbreak{}\texttt{2+}\allowbreak{}\texttt{o(1)},

\[
\begin{split}
 \log_2\left[k_{co}^2g(j)\frac{(ea^2/m_0)^j}{j!}\right]
 &\le2L+\frac{j^2}{2}-j\log_2m_0+O(j\log_2a)\\
 &\le-cL^2.                                              \end{split}
\tag{8.28}
\]

Indeed, the quadratic coefficient is
\texttt{x\^{}2/}\allowbreak{}\texttt{2-x\textless{}}\allowbreak{}\texttt{=}\allowbreak{}\texttt{-3/}\allowbreak{}\texttt{8+}\allowbreak{}\texttt{o(1)}
on this interval. Hence
\texttt{Xi\_}\allowbreak{}\texttt{4=}\allowbreak{}\texttt{2\^{}\{-Omega(L\^{}2)}\allowbreak{}\texttt{\}}.

If instead \texttt{m\_0\textless{}n/N\^{}6}, leave the entire residual
matching unexposed. The near skeleton is a matching, so pointwise

\[
 \beta(M\cup H_{res})\le|E(H_{res})|\le m_0/2,
 \quad
 \sum_e\binom{r_e}{2}\le(a-1)m_0/2.                    \tag{8.29}
\]

All remaining local and topological factors together are therefore at
most
\texttt{exp(O(am\_}\allowbreak{}\texttt{0)}\allowbreak{}\texttt{)}\allowbreak{}\texttt{=}\allowbreak{}\texttt{exp(O(n/}\allowbreak{}\texttt{N\^{}5)}\allowbreak{}\texttt{)}.
In particular, this also upper-bounds the bare middle-cell sum.
Combining (8.16), (8.26), and the two branches gives

\[
 \sum_{\text{all high skeletons}}\text{bare weight}
 \le\exp\{O(\sqrt{nN})+O(N^2)+O(n/N^5)\}
 =\exp\{o(n/N^4)\}.                                     \tag{8.30}
\]

\(\square\)

\section{9. All residual local and cycle
attachments}\label{all-residual-local-and-cycle-attachments}

Fix an actual canonical high skeleton \texttt{(M,j)} as in (8.2), and
let \texttt{m\_0=n-J} be the residual total degree. Its residual degrees
are at most \texttt{U}. Define

\[
 \mathcal A(M,j)=\mathbb E_{res}\left[
  \left(\prod_e g(r'_e)\right)
  2^{\beta(M\cup H_{res})}\mathbf1_{\mathcal E(M,j)}
 \right],                                                \tag{9.1}
\]

where \texttt{E(M,j)} imposes the residual cap
\texttt{r\textquotesingle{}\_e\textless{}=R\_0} and forbids a further
pair in a cell of \texttt{M}. By (6.4), (8.3), and the exact conditional
law, the normalized moment is exactly

\[
 \sum_{(M,j)}\pi(M,j)\left(\prod_{e\in M}g(j_e)\right)
 \mathcal A(M,j).                                       \tag{9.2}
\]

\subsection{Lemma 9.1 (uniform attachment
bound)}\label{lemma-9.1-uniform-attachment-bound}

Uniformly over every canonical high skeleton,

\[
 \ln\mathcal A(M,j)=o(n/N^4).                            \tag{9.3}
\]

More precisely, it is \texttt{O(N\^{}8)} when
\texttt{m\_0\textgreater{}=n/N\^{}6} and \texttt{O(n/N\^{}5)} otherwise.

\subsubsection{Proof: large residual
degree}\label{proof-large-residual-degree}

Suppose \texttt{m\_0\textgreater{}=n/N\^{}6}. Since

\[
 U=(2+o(1))\log_2n,\qquad
 \log_2m_0=\log_2n-O(\log_2N),                           \tag{9.4}
\]

we have \texttt{U\textless{}3\ log\_2\ m\_0} for all sufficiently large
\texttt{n}. For a residual cell outside \texttt{M}, put

\[
 \theta_{ab}=\frac{e d_ad'_b}{m_0},\quad
 \Delta_x=g(x)-g(x-1),                                  \tag{9.5}
\]

and

\[
 \lambda_{ab}=\sum_{x=3}^{R_0}\Delta_x
                    \frac{\theta_{ab}^x}{x!},\qquad
 q_{ab}=\frac{\theta_{ab}^2}{2}+\lambda_{ab}.            \tag{9.6}
\]

Set both quantities to zero on \texttt{M}. Let
\texttt{theta\_*=eU\^{}2/m\_0}. The ratio of consecutive majorants
\texttt{g(x)theta\_*\^{}x/x!} is

\[
 \frac{2^x\theta_*}{x+1},                                \tag{9.7}
\]

and increases with \texttt{x}. Hence the sequence is log-convex. Its
decreasing head is dominated by the \texttt{x=3} term. Consecutive
values of (9.7) grow by the factor
\texttt{2(x+1)/(x+2)\textgreater{}=8/5}, so only \texttt{O(1)}
transition terms can have ratio between \texttt{1/2} and \texttt{1}. Its
increasing tail is dominated by its right endpoint, and

\[
 \log_2\left(g(R_0)\frac{\theta_*^{R_0}}{R_0!}\right)
 \le\frac{R_0^2}{2}-R_0\log_2m_0+O(R_0\log_2U)
 =-\Omega(R_0\log n),                                   \tag{9.8}
\]

using \texttt{R\_0=U/2} and (9.4). Thus, uniformly in the cell,

\[
 \lambda_{ab}\le C\theta_{ab}^3,\qquad
 q_{ab}\le C\theta_{ab}^2.                              \tag{9.9}
\]

Expand each capped local factor as

\[
 g(r'_e)=1+\sum_{x=3}^{R_0}\Delta_x
                   \mathbf1_{\{r'_e\ge x\}}.            \tag{9.10}
\]

If an edge is selected in a fixed even subgraph, use instead

\[
 \mathbf1_{\{r'_e\ge2\}}g(r'_e)
 =\mathbf1_{\{r'_e\ge2\}}
  +\sum_{x=3}^{R_0}\Delta_x\mathbf1_{\{r'_e\ge x\}}.   \tag{9.11}
\]

The terms on the right of (9.11) are alternatives, so a triple cell is
not charged twice. Fix an even subgraph in (6.7), expand (9.10)--(9.11),
and apply Lemma 6.2 jointly to all demanded cells. Dropping the cap only
after this nonnegative expansion gives

\[
 \mathcal A(M,j)
 \le e^{\Lambda_0}
  \sum_{F\text{ even}}\prod_{e\in F\setminus M}q_e,\qquad
 \Lambda_0=\sum_e\lambda_e.                              \tag{9.12}
\]

The degree sums give

\[
 \Lambda_0\le\frac C{m_0^3}
   \left(\sum_a d_a^3\right)\left(\sum_b(d'_b)^3\right)
 \le C\frac{U^4}{m_0},                                  \tag{9.13}
\]

and

\[
 \max_a\sum_bq_{ab}\ \vee\ \max_b\sum_aq_{ab}
 \le C\frac{U^3}{m_0}=:\tau=o(1).                       \tag{9.14}
\]

Every even graph is a union of edge-disjoint simple cycles. Choose one
such decomposition deterministically; dropping edge-disjointness gives

\[
 \sum_{F\text{ even}}\prod_{e\in F\setminus M}q_e
 \le\exp\left\{\sum_{C\text{ simple cycle}}
                  \prod_{e\in C\setminus M}q_e\right\}. \tag{9.15}
\]

For cycles disjoint from \texttt{M}, mark a row start and drop
simplicity and the closing constraint. The weighted row and column sums
are at most \texttt{tau}, and a bipartite cycle has length at least
four. Hence

\[
 \sum_{C:C\cap M=\varnothing}\prod_{e\in C}q_e
 \le\frac{n\tau^4}{1-\tau^2}.                            \tag{9.16}
\]

For a cycle meeting \texttt{M}, mark and orient one matching edge. There
are at most \texttt{2h} choices. Cutting at all matching edges leaves
nonempty residual walks. The residual-walk kernel has row-sum norm at
most

\[
 b=\sum_{\ell\ge1}\tau^\ell=\frac\tau{1-\tau},          \tag{9.17}
\]

while traversal of the matching is a partial permutation of norm one.
Thus, if the relaxed cycle uses \texttt{r} matching edges, its total
remaining weight is at most \texttt{b\^{}r}; there is no fresh factor
\texttt{h} after the first marked edge. Therefore

\[
 \sum_{C:C\cap M\ne\varnothing}
       \prod_{e\in C\setminus M}q_e
 \le2h\sum_{r\ge1}b^r\le Ch\tau.                        \tag{9.18}
\]

Every high cell uses more than \texttt{U/2} stubs, so
\texttt{h\textless{}2n/U}. Combining (9.12)--(9.18),

\[
 \ln\mathcal A(M,j)
 \le C\left\{\frac{U^4}{m_0}+n\tau^4+h\tau\right\}
 =O(N^8)                                                  \tag{9.19}
\]

at the cutoff \texttt{m\_0\textgreater{}=n/N\^{}6}.

\subsubsection{Proof: small residual
degree}\label{proof-small-residual-degree}

Suppose \texttt{m\_0\textless{}n/N\^{}6}. Since \texttt{M} is a
matching, it is a forest, and adding a residual support edge increases
cycle rank by at most one. Hence

\[
 \beta(M\cup H_{res})\le|E(H_{res})|\le m_0/2.           \tag{9.20}
\]

Also, pointwise,

\[
 \sum_e\binom{r'_e}{2}
 \le\frac{U-1}{2}\sum_er'_e=\frac{U-1}{2}m_0.           \tag{9.21}
\]

Equations (6.5), (9.20), and (9.21) bound the integrand in (9.1) by
\texttt{2\^{}\{Um\_0/2\}}. Therefore

\[
 \ln\mathcal A(M,j)=O(Um_0)=O(n/N^5).                   \tag{9.22}
\]

Both (9.19) and (9.22) are \texttt{o(n/N\^{}4)}. \(\square\)

\subsection{Proposition 9.2 (normalized signed second
moment)}\label{proposition-9.2-normalized-signed-second-moment}

For the exact midpoint profile,

\[
 \Lambda_n:=\ln\frac{\mathbb E(Z_{\mathbf k}^{sgn})^2}
                      {(\mathbb EZ_{\mathbf k}^{sgn})^2}
 =o(n/N^4).                                               \tag{9.23}
\]

\subsubsection{Proof}\label{proof-8}

Equation (9.2) is an exact canonical decomposition. Lemma 8.3 sums all
bare high skeletons by \texttt{exp(o(n/N\^{}4))}, including the empty
skeleton, all dense endpoint tables, every near-containment decoration,
and the two middle-strip regimes. Lemma 9.1 bounds the conditional
residual factor by the same exponential scale uniformly over every
skeleton. Multiplication proves the upper bound in (9.23). The
normalized second moment is at least one by variance, so its logarithm
is nonnegative and (9.23) follows. \(\square\)

\section{10. Rare-event amplification}\label{rare-event-amplification}

Proposition 9.2 and (1.5) give the seed

\[
 \Pr\{Z_{\mathbf k}^{sgn}>0\}\ge e^{-\Lambda_n}.        \tag{10.1}
\]

By the observation after (5.19), this implies

\[
 \Pr\{\zeta(G_n)\le k_{co}\}\ge e^{-\Lambda_n}.        \tag{10.2}
\]

We now prove the amplification needed to turn this possibly rare event
into a typical one.

\subsection{Lemma 10.1 (simultaneous leftover
colouring)}\label{lemma-10.1-simultaneous-leftover-colouring}

There is an absolute \texttt{C} such that, with probability
\texttt{1-o(1)}, every \texttt{U\ subseteq{[}n{]}} satisfies

\[
 \chi(G_n[U])\le C\frac{|U|}{N}+n^{1/3}.                 \tag{10.3}
\]

\subsubsection{Proof}\label{proof-9}

Let \texttt{H} be the complement of \texttt{G\_n}. Put
\texttt{u\_0=ceil(n\^{}\{1/4\})}. For any fixed \texttt{u\_0}-set, (1.6)
gives probability
\texttt{exp(-Omega(u\_}\allowbreak{}\texttt{0\^{}2)}\allowbreak{}\texttt{)}
that its \texttt{H}-edge density is below \texttt{1/4}. There are at
most \texttt{exp(u\_0N)} such sets, so a union bound shows that, with
probability \texttt{1-o(1)}, every \texttt{u\_0}-set has density at
least \texttt{1/4}. By averaging over its \texttt{u\_0}-subsets, the
same is then true of every larger set.

Start with any vertex set of size at least \texttt{n\^{}\{1/3\}}. While
its current size is at least \texttt{u\_0}, choose a vertex of maximum
\texttt{H}-degree and replace the current set by its
\texttt{H}-neighbourhood there. Density at least \texttt{1/4} implies
that this neighbourhood has size at least one quarter of the current
size, up to an additive one. The chosen vertices form a clique in
\texttt{H}; before the current size falls from \texttt{n\^{}\{1/3\}} to
\texttt{n\^{}\{1/4\}}, this constructs at least \texttt{cN} vertices for
an absolute \texttt{c\textgreater{}0}. They are an independent set in
\texttt{G\_n}.

For an arbitrary \texttt{U}, repeatedly remove such independent sets and
give each a new colour until fewer than \texttt{n\^{}\{1/3\}} vertices
remain; colour the rest singly. This proves (10.3) simultaneously for
every \texttt{U}. \(\square\)

\subsection{Lemma 10.2 (amplification from a
seed)}\label{lemma-10.2-amplification-from-a-seed}

Suppose deterministic \texttt{k\_n,Lambda\_n\textgreater{}=0} satisfy

\[
 \Pr\{\zeta(G_n)\le k_n\}\ge e^{-\Lambda_n}.             \tag{10.4}
\]

Then for every deterministic \texttt{r\textgreater{}0},

\[
\begin{split}
 \Pr\Bigg\{\zeta(G_n)>k_n+C\bigg(
 &\frac{\sqrt{n\Lambda_n}+\sqrt{nr}}N
   +n^{1/3}+1\bigg)\Bigg\}\\
 &\le e^{-r}+o(1).                                       \end{split}
\tag{10.5}
\]

\subsubsection{Proof}\label{proof-10}

For an induced set \texttt{W}, let
\texttt{zeta(W)}\allowbreak{}\texttt{=}\allowbreak{}\texttt{zeta(G\_}\allowbreak{}\texttt{n{[}W{]}}\allowbreak{}\texttt{)},
and define

\[
 S_k=\max\{|W|:\zeta(W)\le k_n\}.                       \tag{10.6}
\]

Expose the random graph in \texttt{n-1} independent vertex blocks, where
block \texttt{v} contains the edges from \texttt{v} to earlier vertices.
Changing one block changes \texttt{S\_k} by at most one: after deletion
of the affected vertex, every feasible induced set loses at most one
vertex and the two configurations agree on the remainder. Therefore
(1.4) applies.

Since \texttt{S\_k=n} exactly when \texttt{zeta(G\_n)\textless{}=k\_n},
(10.4) and the upper one-sided bounded-differences tail imply

\[
 n-\mathbb ES_k\le\sqrt{(n-1)\Lambda_n/2}.               \tag{10.7}
\]

The lower tail with radius \texttt{sqrt((n-1)r/2)} gives, except with
probability \texttt{e\^{}\{-r\}},

\[
 n-S_k\le
 \sqrt{(n-1)\Lambda_n/2}+\sqrt{(n-1)r/2}.                \tag{10.8}
\]

Choose a maximizing set \texttt{W} and put
\texttt{U={[}n{]}\textbackslash{}setminus\ W}. Concatenating a
\texttt{k\_n}-part cocolouring of \texttt{G{[}W{]}} with an ordinary
colouring of \texttt{G{[}U{]}} gives

\[
 \zeta(G_n)\le k_n+\chi(G_n[U]).                         \tag{10.9}
\]

Intersect (10.8) with the simultaneous event in Lemma 10.1 and apply
(10.3). This proves (10.5). \(\square\)

Apply Lemma 10.2 to (10.2). Put

\[
 B_n=\frac n{N^4},\qquad r_n=\sqrt{B_n}.                 \tag{10.10}
\]

Then \texttt{r\_n-\textgreater{}infinity}, and Proposition 9.2 gives

\[
 \frac{\sqrt{n\Lambda_n}}N=o(n/N^3),\qquad
 \frac{\sqrt{nr_n}}N=o(n/N^3),\qquad
 n^{1/3}=o(n/N^3).                                       \tag{10.11}
\]

Consequently

\[
 \Pr\left\{\zeta(G_n)
 \le k_{co}+o(n/N^3)\right\}\longrightarrow1.          \tag{10.12}
\]

The \texttt{o(n/N\^{}3)} in (10.12) is a deterministic sequence.

\section{11. Completion of the proof}\label{completion-of-the-proof}

Intersect the events (4.5) and (10.12). No independence is required; a
union bound suffices. Equation (5.20) then gives, with high probability,

\[
 \chi(G_n)-\zeta(G_n)
 \ge\left(\frac{q^2\gamma_4}{16}-o(1)\right)
        \frac n{N^3}.                                    \tag{11.1}
\]

For every sufficiently large \texttt{n}, the right side is at least

\[
 c_*\frac n{N^3},\qquad
 c_*:=\frac{q^2\gamma_4}{32}
 =\frac{(\ln2)^2}{32}\ln\frac{200}{153}>0.              \tag{11.2}
\]

This proves (0.1). Since \texttt{n/(ln\ n)\^{}3} tends to infinity, for
every fixed \texttt{M},

\[
 \Pr\{\chi(G_n)-\zeta(G_n)\ge M\}\longrightarrow1.      \tag{11.3}
\]

Every estimate used above is uniform for \texttt{delta\ in\ {[}0,1)} and
for the exact tangent-rounded finite-\texttt{n} profile. The conclusion
therefore holds along all sufficiently large integers, including
sequences approaching either side of every independence-number jump. It
is a high-probability statement, not a density-subset or subsequence
result. \(\blacksquare\)

\section{12. Dependency and provenance
note}\label{dependency-and-provenance-note}

The proof is logically ordered as follows. The elementary phase
expansion and exact first-moment variational calculation construct the
midpoint profile and prove its complete signed first-moment margin. That
margin and elementary independence-set estimates prove the full
partial-diagonal sum. The partial-diagonal sum feeds the dense four-type
transportation estimate. Independently, the exact sign sum and the
configuration model prove the uniform residual-attachment estimate.
These two estimates combine only after both are established, yielding
the normalized second moment. Paley--Zygmund and the proved leftover
amplification are applied last. Thus there is no circular use of a
second-moment conclusion in the profile construction. In particular, no
ordinary-colouring second-moment theorem is invoked at the signed
midpoint, where the ordinary first moment is exponentially small.

For provenance only, the separate dossier files from which the arguments
were consolidated are
\texttt{EXCEPTIONAL\_}\allowbreak{}\texttt{REGIME.md},
\texttt{ALPHA\_}\allowbreak{}\texttt{MINUS\_}\allowbreak{}\texttt{TWO\_}\allowbreak{}\texttt{ROUTE.md},
\texttt{FOUR\_}\allowbreak{}\texttt{SIZE\_}\allowbreak{}\texttt{PARTIAL\_}\allowbreak{}\texttt{RATES.md},
\texttt{DENSE\_}\allowbreak{}\texttt{FOUR\_}\allowbreak{}\texttt{TYPE\_}\allowbreak{}\texttt{MATCHING.md},
\texttt{SIGNED\_}\allowbreak{}\texttt{PROFILE\_}\allowbreak{}\texttt{OVERLAP.md},
\texttt{RESIDUAL\_}\allowbreak{}\texttt{ATTACHMENT.md}, and
\texttt{ALON\_}\allowbreak{}\texttt{CONCENTRATION\_}\allowbreak{}\texttt{EXTENSION.md}.
No step of the proof above depends on those cross-references.

\end{document}
